\documentclass[11pt,a4paper]{article}
\usepackage{fontspec,amsmath,amssymb,amsthm,microtype}
\setmainfont{lmroman10-regular.otf}[BoldFont=lmroman10-bold.otf,ItalicFont=lmroman10-italic.otf,BoldItalicFont=lmroman10-bolditalic.otf]
\usepackage[margin=25mm]{geometry}
\usepackage[unicode,hidelinks]{hyperref}
\hypersetup{pdfauthor={LangevinLean Project}}
\setlength{\emergencystretch}{3em}
\setlength{\parskip}{0.4em}
\setlength{\parindent}{0pt}
\allowdisplaybreaks[2]
\raggedbottom
\title{\texorpdfstring{A sufficient total-\allowbreak{}defect bound with the exact final coefficient}{A sufficient total-defect bound with the exact final coefficient}}
\author{LangevinLean Project}
\date{}
\begin{document}\maketitle
Draft mathematical exposition; these notes may contain mistakes.


The improved rate needs generator cancellation in the cross term, while an ordinary Brownian increment estimate suffices for the total squared local defect. This avoids adding stochastic-integral assumptions to the paper.



For the actual integral solution and its Brownian increment bounds, \(hM\le1\) gives


\[\left\|\int_0^h\bigl(g(Y_t)-g(Y_0)\bigr)\,dt\right\|
\le\frac32Mh\sqrt h\,\sqrt d.\]


The final scalar argument uses \(q=1-mh\), \(m\le A^2\), and the precise implication


\[\begin{gathered}
e_{k+1}^2\le q^2e_k^2+qh^2(MA+\delta)S e_k
+\left(\frac32Mh\sqrt h\,S\right)^2\\
\Longrightarrow\quad e_n\le q^ne_0+\frac{(2MA+\delta)hS}{m}.
\end{gathered}\]

Earlier results used here: \cite[Section~3.11]{LangevinLean2026ULASC}.

\section{\texorpdfstring{Control the total local defect}{Control the total local defect}}


Apply the already proved integral-equation defect bound. A score norm at most \ensuremath{\surd}(Md), Brownian increment norm at most \ensuremath{\surd}(2d)\ensuremath{\surd}t, and \(hM\)\ensuremath{\le}1 give a total defect at most (3/\allowbreak{}2)Mh\ensuremath{\surd}h\ensuremath{\surd}d. The numerical estimate follows from \ensuremath{\surd}2\ensuremath{\le}3/\allowbreak{}2.



Earlier results used here: \cite[Section~3.11.1]{LangevinLean2026ULASC}; \cite[Section~3.11.7]{LangevinLean2026ULASC}.

\section{\texorpdfstring{Preserve the exact order h floor}{Preserve the exact order h floor}}


For a recurrence with cross coefficient (1\ensuremath{-}mh)h\ensuremath{{}^{2}}(MA+\ensuremath{\delta})S and the preceding total defect, use m\ensuremath{\le}A\ensuremath{{}^{2}} to absorb its variance floor. The explicit floors a=h(MA+\ensuremath{\delta})S/\allowbreak{}(2m) and b=3hMAS/\allowbreak{}(2m) sum to at most the manuscript's (2MA+\ensuremath{\delta})\(hS\)/\allowbreak{}m. The final law-level theorem later proves the recurrence rather than assuming it.


\begin{thebibliography}{ULA-LSI}
\bibitem[ULA-SC]{LangevinLean2026ULASC}
LangevinLean Project. \emph{Unadjusted Langevin under Strong Convexity}. LL-01, 2026. Draft report, content revision 6f6f50decdb6. \url{https://odakyildiz.github.io/LangevinLean/docs/pdf/UnadjustedLangevin/FullProof.pdf}.
\end{thebibliography}

\end{document}
